\section{From Labels to a Queue}
\label{sec:queue}

A desk with five agents and a backlog consumes an order, not labels, and two
systems with equal \macrof{} can produce very different orders. We therefore
evaluate the predictions at the level they are used: rather than design a
scoring function, we apply published scheduling rules unmodified and score each
on the cost objective its own theorem is stated over. The simulator is an M/G/5
queue with log-normal service times (mean 8 minutes, $\sigma=0.75$) at
utilisation $\rho$, 200 seeds, with service windows of 30, 120, and 480 minutes
for High, Medium, and Low.

\textbf{The three labels are not three independent signals.} On the English
training track, intent carries 0.7178 nats of mutual information with priority
against priority's entropy of 0.9333 nats, 76.9\% of what there is to know, so
an additive score over an intent head and a priority head adds two estimates of
one quantity and its fitted weight measures collinearity. Sentiment is the
mirror image: 0.0172 nats, of which 86.6\% survives conditioning on intent. It
knows little, and what it knows is not already known. Neither fact is visible in
a weight, which argues for modelling the joint distribution.

Comparing five ways of forming a joint priority posterior, a log opinion pool
reaches 0.2683 log-loss against binary relevance's 0.3697, a 27\% reduction,
while \macrof{} moves across a range of 0.024 with no consistent ordering.
Modelling label dependence improves calibration and not argmax accuracy, which
is exactly the quantity a scheduling index consumes and an accuracy table cannot
report. Classifier chains \cite{read2009classifierchains}, probabilistic chains
\cite{dembczynski2010pcc}, and stacking \cite{wolpert1992stacked} fall between.

\begin{table}[t]
\caption{Published policies at $\rho=1.05$ under the best label model, 200
seeds. Cost columns are mean delay cost, lower is better. Each rule wins on the
objective its own theorem optimizes and is mid-table on the other.}
\label{tab:policies}
\centering
\small
\begin{tabular}{@{}lrrr@{}}
\toprule
Policy & Linear & Convex & High wait (min) \\
\midrule
$c\mu$ / HSF \cite{cox1961queues}        & \textbf{0.121} & 0.137 & 1.77 \\
Generalised $c\mu$ \cite{vanmieghem1995generalized} & 0.149 & \textbf{0.064} & 3.29 \\
Earliest due date                        & 0.146 & 0.111 & 4.78 \\
Accumulating priority                    & 0.197 & 0.099 & 9.34 \\
Static tier (argmax)                     & 0.161 & 0.163 & 6.70 \\
First come, first served                 & 0.381 & 0.625 & 52.46 \\
\bottomrule
\end{tabular}
\end{table}

\textbf{Each rule wins on its own objective.} A bake-off ranking policies on
some other quantity is not testing them. The $c\mu$ rule and Argon and Ziya's
Theorem 3 are stated over linear delay cost and the generalised rule over convex
cost; Table~\ref{tab:policies} measures both. The $c\mu$ rule sweeps linear cost
and sits mid-table on convex, and the generalised rule does the reverse, so the
choice between them is a choice of cost model rather than an empirical contest.

Argon and Ziya's result holds when tested on the objective it is stated over:
against the same label model the posterior index beats the static tier every
time, 0.121--0.125 against 0.161--0.168, ranges that do not overlap, a 25\% cost
reduction. Scored on relative tardiness, which no cited theorem mentions, the
same comparison is a tie and the theorem appears to fail.

\textbf{The ordering rule dominates the probability model.} Under linear cost
the five label models span 3\% while the policies span a factor of $3.1$.
Practitioners choosing where to spend effort should read that ordering. Under
convex cost the label model earns more, spanning 42\%, so dependence modelling
pays once the cost is convex. At $\rho=0.85$ nothing separates: linear cost runs
0.019--0.022 across every combination. With no backlog there is nothing to
order, and every claim here is an overload claim.

\gap{Not in this paper: the cross-script service disparity, which would join
Sections~\ref{sec:script} and \ref{sec:queue} by converting the accuracy gap
into minutes of customer waiting time. The bake-off above scores against gold
priority by design, so classifier error never enters the queue. Recovering it
needs the queue driven by predicted priority with fixed per-track contrasts
against English, plus FIFO and gold-label null controls.}
